% ============================================================
% Section: The Kontablo Agent — Harness Architecture and
%          the Locus of Error
% Placement: after deterministic_logic.tex, before agentic_economy.tex
% Status: new section added Phase 2.5 pre-publication pass
% Epistemic notes:
%   - Claims marked [VERIFICAR] require a literature check before
%     final submission. Do not remove the marker; replace it with
%     a confirmed citation or remove the claim.
%   - Claims marked [PROPUESTA ORIGINAL] are design contributions
%     without prior literature support identified. They must be
%     defended as such in peer review, not presented as established
%     facts.
% ============================================================

\section{The Kontablo Agent: Harness Architecture and the Locus of Error}
\label{sec:harness}

\subsection{The Model--Harness Distinction}
\label{subsec:model_harness}

A common framing of AI reliability in critical systems asks:
\emph{``How do we prevent the LLM from hallucinating?''}
This question, while intuitive, addresses the wrong architectural layer.
It locates the reliability problem inside the model, where the design levers
available to system architects are limited and where the state of the art
offers no near-term guarantees suitable for ledger-grade accuracy.

A more productive framing asks:
\emph{``How do we design the system around the model so that the
decision space the model can influence is constrained to valid outputs,
and anything outside coverage escalates to a human rather than being
confabulated?''} This reframe is the starting point of harness-centric
agent design.

The \textbf{model} is the stochastic component: the Large Language Model that
performs semantic interpretation, natural-language matching, and
open-ended reasoning over unstructured input. The \textbf{harness} is the
deterministic scaffold surrounding it: the tool definitions and schemas
that structure its output, the ontological constraints that bound the valid
output space, the validation layers that reject structurally invalid proposals
before they reach the ledger, and the escalation protocols that route
out-of-coverage cases to human review.

The core practical insight is that the reliability of an agentic system in a
critical domain is determined far more by the harness than by the model's
intrinsic accuracy. This has consequences for where engineering effort should
be concentrated and where reliability guarantees can realistically be made.
Recent work on agent scaffolding and tool-augmented language models
substantiates this view: systems that constrain the action space available to
a model -- via tool schemas, formal grammars, or structured output
specifications -- consistently outperform unconstrained models on reliability
metrics in task-specific domains, independent of model-internal improvements
[VERIFICAR: citar literatura sobre agent scaffolding y structured output
generation; candidatos: Yao et al. 2023 ReAct ICLR; Schick et al. 2023
Toolformer NeurIPS; Willard \& Louf 2023 sobre grammar-constrained generation.
Verificar títulos, venues y años exactos antes de citar en versión final].

\subsection{Ontology-as-Constraint: Eliminating the Classical Accounting
Hallucination Class}
\label{subsec:ontology_constraint}

In the absence of a harness, an LLM deployed for transaction classification
can hallucinate in an accounting-specific way: it may propose an account code
or name that \emph{sounds} plausible but does not correspond to any account in
the client's chart of accounts, or it may classify a transaction to the wrong
node with high expressed confidence and no observable uncertainty signal.

The asymmetry with information retrieval is severe. A 95\% classification
accuracy is acceptable in search; in a general ledger, the 5\% of
misclassified transactions introduces systematic distortions that compound
across reporting periods, survive into audited financial statements, and --
in the presence of the Three-Tier Resolution pipeline -- would do so
\emph{silently}, with no inconsistency flag raised, because an unconstrained
model emitting a plausible-looking but wrong UUID would pass surface validation.

Kontablo addresses this through what we term \textbf{ontology-as-constraint}:
the agent's tool interface (exposed via the MCP server) accepts only Kontablo
Level 3 UUIDs as valid account identifiers. The set of valid UUIDs is finite,
enumerable, and maintained in the ontology graph. The harness rejects any
proposed UUID that does not exist in the graph at validation time, before
the value reaches the ledger write path.

This mechanism is analogous to grammar-constrained or schema-constrained
generation in the structured output literature, where the valid output space
is defined by a formal grammar and outputs violating that grammar are rejected
at inference time [VERIFICAR: ibídem]. In Kontablo's case the ``grammar'' is
the ontology itself: the closed set of valid account UUIDs, each carrying its
nature (debit/credit), IFRS tag, and deterministic boundary predicates from
the Deterministic Boundary Library.

The consequence is that the \emph{classical accounting hallucination class} --
the emission of arbitrary account identifiers that do not exist in the
ontology -- is architecturally eliminated at the harness level.
This is not a probabilistic reduction in hallucination rate; it is a
categorical elimination of a specific error class within the constrained
output space. The model can still be wrong within the set of valid UUIDs
(mapping a transaction to \texttt{expense.admin} when \texttt{expense.cogs}
is correct), but it cannot emit a non-existent account.

\textbf{Implementation dependency:} This guarantee holds only if the harness
enforcement is strict and complete -- no code path allows a UUID to reach the
ledger write step without a graph-membership check. This is a design
requirement, not an automatic consequence of using an LLM with tool
definitions. Any deployment that bypasses the validation layer for performance
or convenience reasons loses this guarantee.

\subsection{The Relocation of Residual Error}
\label{subsec:residual_error}

Eliminating the classical hallucination class does not eliminate error; it
\emph{relocates} it. The residual error population, after ontology-as-constraint
is applied, concentrates at a specific and identifiable location: the
\textbf{boundary of semantic coverage}.

% Figure: Locus of error before and after ontology-as-constraint
% Source: originally in sections/harness_architecture.tex
% Depicts: Three regions (ontology-bounded / coverage boundary / outside coverage)
%          mapped to error disposition (eliminated / relocated / escalated).
%
% Design system: Kontablo/universal=blue!10, boundary/warning=yellow!20,
%   outside=gray!15, eliminated=green!12, residual-error=red!10.
%   Arrows >=Latex, thick. rounded corners=4pt.
\begin{figure}[ht]
\centering
\fitwidth{%
\begin{tikzpicture}[
    >=Latex,
    node distance=1.5cm,
    regionnode/.style={draw=blue!50, rectangle, rounded corners=4pt,
        minimum width=6.0cm, minimum height=1.1cm,
        align=center, font=\small},
    elimnode/.style={draw=green!60!black, rectangle, rounded corners=4pt,
        fill=green!12, text width=5.2cm, align=center,
        minimum height=1.0cm, font=\small},
    errornode/.style={draw=red!60!black, rectangle, rounded corners=4pt,
        fill=red!10, text width=5.2cm, align=center,
        minimum height=1.0cm, font=\small},
    edgelabel/.style={font=\scriptsize, fill=white, inner sep=1pt}
]
\node[regionnode, fill=blue!10] (bounded)
    {Ontology-bounded decision space\\(30 Level~3 UUIDs)};
\node[regionnode, fill=yellow!15, draw=yellow!70!black, below=0.8cm of bounded] (boundary)
    {Semantic coverage boundary\\(unmapped types, jurisdictional gaps,
    ambiguous cases)};
\node[regionnode, fill=gray!15, draw=gray!55, below=0.8cm of boundary] (outside)
    {Outside ontology coverage\\(routed to human via CRA)};

\node[elimnode, right=2.0cm of bounded]  (e1)
    {Classical hallucination class\\\textbf{eliminated} by harness};
\node[errornode, right=2.0cm of boundary] (e2)
    {Residual error population\\\textbf{relocated here}};
\node[elimnode, right=2.0cm of outside]  (e3)
    {No silent error:\\escalated to human co-signer};

\draw[->, thick] (bounded)  -- (e1);
\draw[->, thick] (boundary) -- (e2);
\draw[->, thick] (outside)  -- (e3);
\end{tikzpicture}%
}
\caption{Locus of error before and after ontology-as-constraint is applied.
Classical hallucination (arbitrary invalid outputs) is eliminated at the
harness layer. Residual errors relocate to the semantic coverage boundary:
unmapped transaction types, incomplete jurisdictional rules, and semantically
ambiguous cases within valid UUID space. Cases outside coverage are escalated
to human review rather than silently committed.}
\label{fig:locus_of_error}
\end{figure}


Errors at the semantic coverage boundary take three forms:

\begin{enumerate}
    \item \textbf{Unmapped transaction types.} A transaction type not yet
    represented in the ontology has no valid UUID. The Three-Tier Resolution
    Pipeline routes it to Tier 3 (semantic AI fallback); if confidence remains
    below the threshold (Section~6), it is flagged for human review rather than
    forcing a low-confidence posting. The error is not confabulated -- it is
    escalated and made observable.

    \item \textbf{Jurisdictional edge cases.} A jurisdictional rule not yet
    encoded in the Deterministic Boundary Library may allow a semantically
    plausible UUID to pass boundary validation while violating an implicit
    local constraint. These errors are attributable to incomplete rule coverage,
    not to model hallucination. They are also correctable through ordinary
    ontology maintenance as cases are discovered.

    \item \textbf{Semantic ambiguity within coverage.} Some transactions
    satisfy the boundary predicates of more than one Level 3 node. In these
    cases Tier 3 may produce a defensible but suboptimal mapping. The
    Co-responsibility Architecture (Section~10) manages this residual by
    requiring human review for low-confidence Tier 3 outputs.
\end{enumerate}

This relocation is architecturally significant because boundary errors are
\textbf{observable}: they manifest as flagged inconsistencies, low-confidence
outputs, and unmapped-transaction queues -- all instrumentable signals.
Stochastic interior errors from an unconstrained model are, by contrast,
\textbf{silent}: a confident-but-wrong classification produces no observable
artifact beyond the wrong posting itself, invisible to the audit trail unless
a human catches it retrospectively.

\textbf{[PROPUESTA ORIGINAL: La formulación de ``relocalización del error''
como categoría analítica distinta de ``reducción de la tasa de error'' no
tiene, a nuestro conocimiento, correlato directo en la literatura existente
sobre agentes LLM en sistemas críticos. Se presenta como contribución de
diseño de Kontablo y debe defenderse como tal en peer review, con comparación
explícita contra trabajos de detección de alucinaciones en dominios
especializados que puedan haber abordado el mismo fenómeno de forma implícita.
Realizar búsqueda bibliográfica antes de la versión final.]}

\subsection{The Engineering Reframe: From Intractable to Tractable}
\label{subsec:reframe}

The classical framing -- ``how do we stop the model from hallucinating?'' --
is not only intractable as stated, but directs engineering effort at the wrong
target. Hallucination in the general sense is an active area of model research
with no near-term resolution that guarantees the classification accuracy a
ledger requires across all inputs and distribution shifts.

The productive reframe, which Kontablo's architecture instantiates, is:

\begin{quote}
\emph{How do we design the harness so that (a) the output space the agent
can emit is bounded by the ontology, making confabulation of non-existent
accounts architecturally impossible; (b) decisions within the bounded space
that can be resolved by deterministic logic are resolved without inference;
and (c) cases outside coverage or below confidence threshold are escalated
to a human co-signer rather than silently committed?}
\end{quote}

This is a tractable engineering problem with identifiable components:

\begin{itemize}
    \item Schema-constrained tool interfaces that enforce UUID membership
    \item The Three-Tier Resolution Pipeline (Section~6), which maximizes
    the proportion of decisions resolved deterministically before reaching
    the model
    \item The Deterministic Boundary Library (Section~6.3), which validates
    even Tier 3 proposals against ontological predicates before ledger write
    \item The Co-responsibility Architecture (Section~10), which governs
    the out-of-coverage escalation path
\end{itemize}

Kontablo's architectural principle~5 -- ``determinism over stochasticity
wherever movable'' (ADR-009) -- formalizes this reframe as a standing design
driver: whenever a decision can be migrated from stochastic inference to
deterministic logic, it must be. The harness is the system-level
implementation of this principle. The ontology-as-constraint is its primary
instrument.

\subsection{The Co-responsibility Architecture as Residual Error
Management}
\label{subsec:cra_as_error_management}

The Co-responsibility Architecture (ADR-008, Section~10) is commonly
presented as a governance mechanism for human oversight. In the context of
the harness architecture, it plays a more specific functional role: it is the
\textbf{residual error management layer}.

Where the three-tier pipeline and the Deterministic Boundary Library resolve
decisions within coverage, the CRA governs what happens at and beyond the
coverage boundary. When the harness cannot deterministically resolve a
transaction -- because no Tier 1 or Tier 2 rule matches, and Tier 3 returns
confidence below threshold -- the CRA routes the transaction to human review,
enforces justification logging, and records the outcome with an immutable
audit trail.

The human accountant in this design is not a rubber-stamp approver of AI
output. They are the \textbf{resolver of the residual error population}: the
cases that are genuinely beyond automated coverage and require human judgment.
The CRA's inconsistency flagging mechanism ensures that every such resolution
is recorded and attributable, operationalizing the ISACA shared-responsibility
model and the EU AI Act's requirements for high-risk AI systems in financial
contexts (ISACA, 2024; EU AI Act, 2025).

This framing implies a property of the system over time: as the ontology
expands -- more jurisdictions mapped, more transaction types covered, more
boundary predicates encoded -- the population of out-of-coverage cases
shrinks and the human review load decreases. The CRA, operating on a
shrinking residual, becomes progressively less burdened without ever being
removed. The coverage boundary never reaches zero; there will always be novel
transactions, regulatory changes, and jurisdictional edge cases that require
human judgment. The design accepts this reality and manages it, rather than
claiming to eliminate it.

\subsection{Future Extension: Per-Client Determinism (ADR-011)}
\label{subsec:adr011_preview}

The harness principle admits a natural extension to the per-client layer.
A Client-Specific Determinism Agent (proposed in ADR-011) would observe an
individual client's transaction history, identify recurring patterns currently
resolved via Tier 3 inference, and propose client-specific Tier 1 or Tier 2
rules for those patterns. This extends the harness locally, without modifying
the shared ontology.

The same co-responsibility constraint applies: no client-specific rule is
activated without explicit human co-signature from the client's designated
accountant. The same data sovereignty constraint applies: the agent operates
entirely on client infrastructure.

ADR-011 is explicitly deferred to Phase~3+ and is documented here to
establish its architectural lineage: it is an application of the same harness
design principle at a finer (per-client) granularity, not a new concept.
The iterative loop of ADR-009 -- continuously identifying decisions that can
be migrated from inference to deterministic logic -- extends naturally and
explicitly to the client layer once production deployments provide the
empirical base required.
